\section*{Bab \ref{ch:g10:electron-shells} --- Kulit Elektron dan Tabel Periodik}

\begin{solution}{exo:g10:electron-shells:1}
C: $1\mathrm{s}^2\,2\mathrm{s}^2\,2\mathrm{p}^2$. N:
$1\mathrm{s}^2\,2\mathrm{s}^2\,2\mathrm{p}^3$. Mg:
$1\mathrm{s}^2\,2\mathrm{s}^2\,2\mathrm{p}^6\,3\mathrm{s}^2$.
\end{solution}

\begin{solution}{exo:g10:electron-shells:2}
Karbon 4, nitrogen 5, magnesium 2.
\end{solution}

\begin{solution}{exo:g10:electron-shells:3}
\omterm{def:g10:electron-shells:configuration}{Kulit terluar} $n = 3$: periode 3. $2 + 3 = 5$ \omterm{def:g10:electron-shells:valence}{elektron valensi}: \omterm{def:g9:periodic-table-first-look:period}{golongan}
15 (fosfor).
\end{solution}

\begin{solution}{exo:g10:electron-shells:4}
\omterm{def:g10:electron-shells:family}{Logam alkali} (litium, natrium, kalium); \omterm{def:g10:electron-shells:family}{halogen} (fluorin, klorin,
bromin, iodin); \omterm{def:g10:electron-shells:family}{gas mulia} (helium, neon, argon).
\end{solution}

\begin{solution}{exo:g10:electron-shells:5}
\omterm{def:g7:atoms-and-molecules:atom}{Atom-atom} cenderung menerima atau melepaskan \omterm{def:g9:inside-the-atom:nucleus}{elektron} sehingga memiliki
delapan \omterm{def:g9:inside-the-atom:nucleus}{elektron} di \omterm{def:g10:electron-shells:configuration}{kulit terluarnya}, seperti \omterm{def:g10:electron-shells:family}{gas mulia} yang terdekat.
\end{solution}

\begin{solution}{exo:g10:electron-shells:6}
K, $2.8.8.1$, melepaskan satu: \ce{K+},
$1\mathrm{s}^2\,2\mathrm{s}^2\,2\mathrm{p}^6\,3\mathrm{s}^2\,3\mathrm{p}^6$
(argon). F, $2.7$, menerima satu: \ce{F-},
$1\mathrm{s}^2\,2\mathrm{s}^2\,2\mathrm{p}^6$ (neon).
\end{solution}

\begin{solution}{exo:g10:electron-shells:7}
$1\mathrm{s}^2\,2\mathrm{s}^2\,2\mathrm{p}^6\,3\mathrm{s}^2$: $Z = 12$,
magnesium.
\end{solution}

\begin{solution}{exo:g10:electron-shells:8}
\ce{Mg^{2+}} dan \ce{Cl-}: \ce{MgCl2}. \ce{Al^{3+}} dan \ce{O^{2-}}:
\ce{Al2O3}.
\end{solution}

\begin{solution}{exo:g10:electron-shells:9}
Belerang (6 \omterm{def:g10:electron-shells:valence}{elektron valensi}, $2.8.6$), dalam \omterm{def:g9:periodic-table-first-look:period}{golongan} yang sama, 16.
\end{solution}

\begin{solution}{exo:g10:electron-shells:10}
\omterm{def:g10:electron-shells:configuration}{Kulit terluarnya} sudah penuh (delapan \omterm{def:g9:inside-the-atom:nucleus}{elektron}): neon tidak memperoleh
apa-apa dengan melepaskan atau menerima \omterm{def:g9:inside-the-atom:nucleus}{elektron}.
\end{solution}

\begin{solution}{exo:g10:electron-shells:11}
Argon memiliki 18 \omterm{def:g9:inside-the-atom:nucleus}{elektron}; \ce{X^{2+}} telah melepaskan 2, jadi X
memiliki 20: kalsium.
\end{solution}

\begin{solution}{exo:g10:electron-shells:12}
Helium, $1\mathrm{s}^2$, memiliki \omterm{def:g10:electron-shells:configuration}{kulit terluar} yang penuh (\omterm{def:g10:electron-shells:configuration}{kulit 1} hanya
menampung dua \omterm{def:g9:inside-the-atom:nucleus}{elektron}), seperti \omterm{def:g10:electron-shells:family}{gas mulia} lainnya; helium tidak
berperilaku seperti \omterm{def:g9:periodic-table-first-look:metal}{logam-logam} \omterm{def:g9:periodic-table-first-look:period}{golongan} 2.
\end{solution}

\begin{solution}{exo:g10:electron-shells:13}
\omterm{def:g10:electron-shells:valence}{Elektron valensinya} berada di kulit 4, lebih jauh dari inti daripada
\omterm{def:g10:electron-shells:valence}{elektron valensi} natrium di \omterm{def:g10:electron-shells:configuration}{kulit 3}: karena kurang kuat terikat,
\omterm{def:g9:inside-the-atom:nucleus}{elektron} itu lebih mudah lepas.
\end{solution}

\begin{solution}{exo:g10:electron-shells:14}
Keempatnya memiliki 18 \omterm{def:g9:inside-the-atom:nucleus}{elektron} dan konfigurasi argon:
\[ 1\mathrm{s}^2\,2\mathrm{s}^2\,2\mathrm{p}^6\,3\mathrm{s}^2\,3\mathrm{p}^6. \]
\end{solution}

\begin{solution}{exo:g10:electron-shells:15}
Unsur itu menerima 3 \omterm{def:g9:inside-the-atom:nucleus}{elektron} untuk mencapai oktet: unsur itu memiliki 5
\omterm{def:g10:electron-shells:valence}{elektron valensi}, \omterm{def:g9:periodic-table-first-look:period}{golongan} 15. Pada periode 3, unsur itu adalah fosfor:
\[ 1\mathrm{s}^2\,2\mathrm{s}^2\,2\mathrm{p}^6\,3\mathrm{s}^2\,3\mathrm{p}^3. \]
\end{solution}

\begin{solution}{pb:g10:electron-shells:1}
\textbf{1.} Aluminium: 13 \omterm{def:g9:inside-the-atom:proton}{proton}, 13 \omterm{def:g9:inside-the-atom:nucleus}{elektron}. Oksigen: 8 \omterm{def:g9:inside-the-atom:proton}{proton}, 8
\omterm{def:g9:inside-the-atom:nucleus}{elektron}.

\textbf{2.} Al: $1\mathrm{s}^2\,2\mathrm{s}^2\,2\mathrm{p}^6\,3\mathrm{s}^2\,3\mathrm{p}^1$.
O: $1\mathrm{s}^2\,2\mathrm{s}^2\,2\mathrm{p}^4$.

\textbf{3.} Al: 3 \omterm{def:g10:electron-shells:valence}{elektron valensi}, periode 3, \omterm{def:g9:periodic-table-first-look:period}{golongan} 13. O: 6,
periode 2, \omterm{def:g9:periodic-table-first-look:period}{golongan} 16.

\textbf{4.} Aluminium adalah \omterm{def:g9:periodic-table-first-look:metal}{logam}, oksigen \omterm{def:g9:periodic-table-first-look:metal}{nonlogam}.

\textbf{5.} \ce{Al^{3+}}: $1\mathrm{s}^2\,2\mathrm{s}^2\,2\mathrm{p}^6$.

\textbf{6.} \ce{O^{2-}}: $1\mathrm{s}^2\,2\mathrm{s}^2\,2\mathrm{p}^6$.

\textbf{7.} Neon.

\textbf{8.} Aluminium melepaskan 3, oksigen menerima 2.

\textbf{9.} $2 \times (+3) + 3 \times (-2) = 0$: \ce{Al2O3}.

\textbf{10.} Dilepaskan: $2 \times 3 = 6$. Diterima: $3 \times 2 = 6$.

\textbf{11.} \omterm{def:g9:inside-the-atom:nucleus}{Elektron} tidak diciptakan atau dimusnahkan: \omterm{def:g9:inside-the-atom:nucleus}{elektron} yang
dilepaskan aluminium tepat sama dengan yang diterima oksigen, dan
senyawanya netral.

\textbf{12.} \omterm{def:g9:ions:ion}{Ion} itu bermuatan sama, $+3$: muatan-muatan dalam kristal
tetap seimbang.

\textbf{13.} $2 \times 27 + 3 \times 16 = 102$ \omterm{def:g9:inside-the-atom:proton}{nukleon}.

\textbf{14.} $102 \times 1.67 \times 10^{-27} = \qty{1.70e-25}{kg}$.

\textbf{15.} $10^{-3} / 1.70 \times 10^{-25} \approx 5.87 \times 10^{21}$
satuan rumus.

\textbf{16.} $2 \times 5.87 \times 10^{21} = 1.17 \times 10^{22}$
\omterm{def:g9:ions:ion}{ion} \ce{Al^{3+}}; $3 \times 5.87 \times 10^{21} = 1.76 \times 10^{22}$
\omterm{def:g9:ions:ion}{ion} \ce{O^{2-}}.

\textbf{17.} Sebuah \omterm{def:g9:inside-the-atom:nucleus}{elektron} kira-kira 1836 kali lebih ringan daripada
sebuah \omterm{def:g9:inside-the-atom:proton}{nukleon}; ke-50 \omterm{def:g9:inside-the-atom:nucleus}{elektron} dalam satu satuan rumus beratnya kurang
dari tiga perseratus satu \omterm{def:g9:inside-the-atom:proton}{nukleon}.

\textbf{18.} $6 \times 5.87 \times 10^{21} \approx 3.5 \times 10^{22}$
\omterm{def:g9:inside-the-atom:nucleus}{elektron} berpindah untuk \qty{1}{g} korundum.
\end{solution}
